\documentclass[]{../../../tex/classes/homework}
\usepackage{../../../tex/packages/hebrewSupport}
\usepackage{../../../../cs-tau-shpe/tex/packages/mathShortcuts}

\usepackage{pgfplots}
\let\mrm \mathrm

\author{שחר פרץ}
\title{שיעורי בית בפיזיקה}
\begin{document}
	\maketitle
	\section{}
	\begin{enumerate}[A.]
		\item מהירות ממוצעת הינה $\bar v := \frac{\Delta x}{\Delta t}$ על פני פרק זמן מסוים. 
		\item המהירויות המוצעות של סירה א' ושל סירה ב' זהות, שכן הפרס המרחקים זהה ($4\mathrm{km}$) והפרש הזמנים זהה ($230\mathrm{s}$). 
		\item תאוצת סירה ב' שלילית, שכן הפרבולה עצובה (המהירות פוחתת לאורך זמן) דהיינו המקדם של $x^{2}$ שלילי. 
		\item נמצא את גודל תאוצתה של סירה ב' החל מ־$t = 150\mrm{s}$. נתחיל מלחשב את המהירות ההתחלתית: 
		\[ v_{150} = \frac{3\mrm{km}}{150\mrm{s}} = 20\mrm{\frac{m}{s}} \]
		משוואת הזמן־מיקום: 
		\[ x(\Delta t) = x_{150} + v_{150}\Delta t + a\Delta t^{2} \]
		ומכאן ידוע: 
		\[ x(80) = x(230 - 150) = 4\mrm{km} \implies 3000 + 80 \cdot 20 + a \cdot 80^{2} = 4000 \implies 3000 + 1600 + a \cdot 6400 = 4000 \]
		נמשיך: 
		\[ \implies a \cdot 6400 = -600 \implies a = -\frac{3}{32}\mrm{\frac{m}{s^{2}}} \]
		\item להלן גרף המהירות־זמן של המכונית: 
		\sen
		\begin{center}
			\begin{tikzpicture}
				\begin{axis}[axis lines=center, ymin=0, xmax=250, xtick={0, 50, 100, 150, 200, 230, 250}, ytick={0, 5, 10, 15, 20}, xlabel=$t\csb{\mrm{s}}$, ylabel=$v\csb{\mrm{\frac{m}{s}}}$]
	    			\addplot [domain=0:150] { 20 };
	    			\addplot [domain=150:230] { -3/32*(x-150) + 20 };
				\end{axis}
			\end{tikzpicture}
		\end{center}
		\she
	\end{enumerate}
	
	\section{}
	בה''כ ב־$t = 0$ מיקום המכוניות מיקום $0$. 
	\begin{enumerate}[A.]
		\item גרף א' שייך למכונית A בעוד גרף ב' שייך למכונית B, שכן תאוצה שלילית משמעה שיפוע שלילי של הגרף בסביבת $t = 0$, ותאוצה חיובית משמעה שיפוע חיובי של הגרף בסביבת $t = 0$ (הנקובה). 
		\item[ב+ג.] ננסה להבין מתי המכוניות חולפות זו לצד זו בפעם השנייה. לשם כך, ננסה להבין כמה הם התרחקו במקטע הראשון: 
		\[ \begin{cases}
			a_A^{1} = \frac{12 - 4}{2 - 0} = \frac{8}{2} = 4\mrm{\frac{m}{s^{2}}} \\
			a_B^{1} = \frac{6 - 10}{2 - 0} = \frac{-4}{2} = -2\mrm{\frac{m}{s^{2}}}
		\end{cases} \quad \begin{cases}
			a_A^{2} = \frac{0 - 12}{4 - 2} = \frac{-12}{2} = -6\mrm{\frac{m}{s^{2}}} \\
			a_B^{2} = \frac{18 - 6}{4 - 2} = \frac{12}{2} = 6\mrm{\frac{m}{s^{2}}}
		\end{cases} \]
		\[ \begin{cases}
			x_A^{1}(t) = 0 + 4t + 4t^{2} &\ \dequad \implies x_A^{1}(2) = 8 + 16 = 24\mrm{m} \\
			x_B^{1}(t) = 0 + 10t - 2t^{2} &\ \dequad \implies x_B^{A}(2) = 20 - 4 = 16\mrm{m}
		\end{cases} \]
		נבדוק האם הן נפרשות במקטע זה: 
		\[ x_A^{1}(t) = x_B^{1}(t) \iff 4t + 4t^{2} = 10t - 2t^{2} \iff 6t^{2} = 6t \iff t = 1\mrm{s} \]
		דהיינו הן נפרשות לראשונה לאחר שנייה. 
		
		מכאן שבמקטע השני: 
		\[ \begin{cases}
			x_A^{2}(\Delta t) = 24 + 12t - 6t^{2} \\
			x_B^{2}(\Delta t) = 16 + 6t + 6t^{2}
		\end{cases} \]
		נחסר את המשוואות ע''מ למצוא מתי המכוניות יפגשו במקטע השני: 
		\[ x_A^{2}(\Delta t) = x_B^{2}(\Delta t) \iff 24 +12t  - 6t^{2} = 16 + 6t + 6t^{2} \iff 8 + 6t - 12t^{2} = 0 \]
		הדיסקרימיננטה חיובית, ומכאן שהמכוניות יפגשו בהמשך בזמנים: 
		\[ \Delta t = \frac{1}{4}\pm 0.854 \implies t = 2.25 \pm 0.854 \mrm{s} \]
		נבחין שהפתרון $2.25 - 0.854$ אינו בתחום ההגדרה $t \in [2, 4]$ למקטע השני, ומכאן שאיננו תקף. לכן הן אך ורק יפגשו בפעם השנייה ב־$t = 3.104\mrm{s}$, ללא פגישה בפעם השלישית. 
		\item התלמיד הטוען שהמהירויות הממוצעות של המכוניות שונו, טועה: ההעתק שהמכוניות עושות ב־$[0, 2]$ שונה (כפי שחישבנו בסעיף הקודם), ומשום ש־$\Delta t = 2$ קבוע, המהירויות הממוצעות שלהן שונות. מכאן שתלמיד א' צודק. 
		\item ראשית מכונית ב' מתחילה במהירות גדולה יותר ומקדימה את מכונית א במקטע $(0, 1)$. ב־$t = 1$ הן נפגשות ובמקטע $(1, 3.104)$ מכונית א' מובילה. אך ב־$t = 3.104$ הן נפגשות בפעם האחרונה ומכאן ואילך במקטע $(3.104, \inft)$ מכונית ב' מובילה. 
	\end{enumerate}
	
	
	\ndoc
\end{document}